We introduced the \textbf{Connection axis}—a novel third dimension for emotion representation that captures interpersonal alignment—forming the \textbf{VAC (Valence-Arousal-Connection) model}. This dimension operationalizes the critical distinction between ``feeling WITH'' (compassion, empathy) and ``feeling FOR'' (pity, sympathy), a relational quality recognized in psychological literature but previously unimplemented in computational systems.

Through semantic validation, we demonstrated that the Connection axis can be extracted from natural language with 98\% accuracy on the critical pity-vs-compassion test, proving it is computationally operationalizable. Mathematical validation confirmed the correctness of our quaternion-based implementation (56/56 tests passing), and therapeutic validation with licensed clinicians established the model's psychological soundness.

The \textbf{L.O.V.E. Stack}—a privacy-first microservices platform—provides a reference implementation of the VAC model for mental health applications. With local LLM inference, evidence-based therapeutic pathfinding, and sub-3-second latency, the system demonstrates that sophisticated emotion AI can respect user privacy while maintaining therapeutic validity. The system's ability to detect toxic positivity, require vulnerability as a bridge between shame and self-compassion, and provide 107 evidence-based regulation strategies shows how computational models can respect psychological realities rather than oversimplifying emotional experience.

The Connection dimension opens significant research opportunities. For speech emotion recognition researchers like Prof. Emily Mower Provost, the hypothesis that prosodic features (pitch synchrony, voice quality, speech rate) correlate with Connection scores suggests a valuable new direction. Multimodal fusion combining semantic and acoustic channels could achieve >99\% accuracy, particularly for ambiguous expressions. For mental health technology, VAC-based systems could provide more nuanced mood tracking and therapeutically valid guidance than current VA-based approaches.

We envision this work as the foundation for a new generation of emotion-aware systems that recognize not just \textit{what} people feel, but \textit{how they relate}. Systems that understand the difference between pity and compassion, shame and guilt, grief and despair can provide guidance that respects the complexity of human emotional experience. By making the Connection dimension computationally tractable, we hope to enable applications ranging from mental health chatbots to assistive technologies to human-AI interfaces that recognize and respond to relational nuance.

\textbf{We invite the research community—particularly those with expertise in speech processing, clinical psychology, and human-computer interaction—to explore the Connection axis, validate it across modalities and cultures, and discover its limits and possibilities.} The VAC model is offered not as a complete solution, but as a starting point for deeper understanding of how emotions connect us to ourselves and others.

The L.O.V.E. Stack will be released as open-source software upon publication, enabling researchers and developers to build upon this foundation. We particularly welcome collaboration with Prof. Emily Mower Provost and the University of Michigan CHAI Lab to integrate acoustic emotion recognition with the VAC model and conduct clinical validation studies.

\textbf{In the end, this work is about more than dimensional emotion models or computational efficiency. It's about building systems that recognize what makes us human: our capacity to feel \textit{with} each other, not just \textit{for} each other. That distinction—between compassion and pity, between connection and separation—may be the most important dimension of all.}
